\chapter{Wick Theorem and Diagrammatics}

\begin{remarkbox}
The Dyson series expresses interacting quantities as time-ordered products of
free fields. Wick's theorem is the algebraic tool that evaluates those products.
It reduces them to normal-ordered products plus contractions, and each
contraction is a propagator.
\end{remarkbox}

\section{Why Wick's Theorem Appears}

In Chapter 7, perturbation theory led to expressions of the form
\[
    \langle0|T\{\hat\phi_I(x_1)\cdots\hat\phi_I(x_n)\}|0\rangle.
\]
The fields are free interaction-picture fields, but there may be many of them.
To compute perturbative corrections, one needs a systematic way to evaluate such
vacuum expectation values.

Wick's theorem provides exactly this. It says that a time-ordered product of
free fields can be rewritten as a sum of normal-ordered products multiplied by
contractions. In vacuum expectation values, the normal-ordered parts vanish
unless all fields are contracted. Thus vacuum expectation values reduce to sums
of products of propagators.

This is the algebraic origin of Feynman diagrams.

\section{Normal Ordering Revisited}

Normal ordering places all creation operators to the left of all annihilation
operators. For a scalar field, we write
\[
    :\hat a_{\vect{k}}\hat a_{\vect{p}}^\dagger:
    = \hat a_{\vect{p}}^\dagger\hat a_{\vect{k}}.
\]
For fields, normal ordering means applying this rule after expanding the field
into creation and annihilation parts. Schematically,
\[
    \hat\phi(x)=\hat\phi^{(+)}(x)+\hat\phi^{(-)}(x),
\]
where \(\hat\phi^{(+)}\) contains annihilation operators and
\(\hat\phi^{(-)}\) contains creation operators.

A normal-ordered product has zero vacuum expectation value unless it is just a
number:
\[
    \langle0|:\hat\phi(x_1)\cdots\hat\phi(x_n):|0\rangle=0
\]
for \(n>0\). This is why normal ordering is useful in perturbation theory.

\section{Contractions}

For two scalar fields, the contraction is defined as the difference between the
time-ordered product and the normal-ordered product:
\[
    \operatorname{contr}\!\bigl(\hat\phi(x),\hat\phi(y)\bigr)
    \equiv
    T\{\hat\phi(x)\hat\phi(y)\}
    - :\hat\phi(x)\hat\phi(y):.
\]
For a free real scalar field, this contraction is the Feynman propagator:
\[
    \operatorname{contr}\!\bigl(\hat\phi(x),\hat\phi(y)\bigr)
    = D_F(x-y).
\]
Equivalently,
\[
    \langle0|T\{\hat\phi(x)\hat\phi(y)\}|0\rangle
    = D_F(x-y).
\]
A contraction is therefore not a new physical object. It is the propagator
written as an algebraic operation inside a time-ordered product.

\section{Wick Theorem for Scalar Fields}

For free scalar fields, Wick's theorem states that
\[
    T\{\hat\phi_1\hat\phi_2\cdots\hat\phi_n\}
    = :\hat\phi_1\hat\phi_2\cdots\hat\phi_n:
      + \text{all single contractions}
      + \text{all double contractions}
      + \cdots,
\]
where \(\hat\phi_i=\hat\phi(x_i)\). The expansion continues until all possible
contractions have been included.

For two fields,
\[
    T\{\hat\phi_1\hat\phi_2\}
    = :\hat\phi_1\hat\phi_2:
      + D_F(x_1-x_2).
\]
For four fields,
\[
\begin{aligned}
    \langle0|T\{\hat\phi_1\hat\phi_2\hat\phi_3\hat\phi_4\}|0\rangle
    ={}& D_F(x_1-x_2)D_F(x_3-x_4)\\
       &+D_F(x_1-x_3)D_F(x_2-x_4)\\
       &+D_F(x_1-x_4)D_F(x_2-x_3).
\end{aligned}
\]
Only fully contracted terms survive between vacuum states.

\begin{examplebox}
The vacuum expectation value of an odd number of free real scalar fields
vanishes. There is no way to pair all fields into contractions.
\end{examplebox}

\section{Application to the Dyson Series}

Consider \(\phi^4\) theory with
\[
    \mathcal{H}_{\mathrm{int}}(x)
    = \frac{\lambda}{4!}\hat\phi_I(x)^4.
\]
The first-order correction to a two-point function contains
\[
    \frac{-i\lambda}{4!}\int\dd^4z\,
    \langle0|T\{\hat\phi(x)\hat\phi(y)\hat\phi(z)^4\}|0\rangle.
\]
Wick's theorem reduces this six-field expectation value to a sum of products of
three propagators. Some contractions connect \(x\) and \(y\) to the interaction
point \(z\); others form vacuum bubbles or tadpole factors.

This example illustrates the general rule: interaction terms add field
insertions, and Wick's theorem pairs those insertions into propagators.

\section{Wick Theorem for Fermions}

For fermionic fields, Wick's theorem has the same structure but includes signs.
The basic fermion contraction is
\[
    \operatorname{contr}\!\bigl(\hat\psi(x),\hat{\bar\psi}(y)\bigr)
    = S_F(x-y),
\]
where \(S_F\) is the Dirac propagator. Time ordering of fermions includes a
minus sign when two fermionic operators are exchanged.

For example,
\[
    T\{\hat\psi(x)\hat{\bar\psi}(y)\}
    = \theta(x^0-y^0)\hat\psi(x)\hat{\bar\psi}(y)
      -\theta(y^0-x^0)\hat{\bar\psi}(y)\hat\psi(x).
\]
The minus sign is not optional. It is required by the anticommutation relations
of fermionic fields.

In diagrammatic language, every closed fermion loop contributes an additional
minus sign. This sign has its origin in the permutations required to contract
fermionic operators.

\section{Fermionic Sign Factors}

When applying Wick's theorem to fermions, one must count how many exchanges of
fermionic operators are needed to bring the contracted fields next to each
other. Each exchange contributes a minus sign.

A useful rule is:
\[
    \text{fermionic sign} = (-1)^{N_{\mathrm{exchanges}}}.
\]
For practical Feynman rules, these signs are packaged into standard rules:
closed fermion loops give a minus sign, and external fermion lines must be
ordered consistently with the chosen amplitude convention.

The sign bookkeeping is one of the main reasons diagrammatic notation is useful.
It turns long operator manipulations into a structured set of graphical and
algebraic rules.

\section{From Algebra to Diagrams}

Each term produced by Wick's theorem can be represented graphically. The rules
are:
\begin{itemize}
    \item each field insertion is represented by a point or an external leg;
    \item each contraction is represented by a line;
    \item each interaction Hamiltonian insertion is represented by a vertex;
    \item each vertex is integrated over spacetime;
    \item each line contributes the propagator of the corresponding field.
\end{itemize}
For \(\phi^4\) theory, each interaction vertex has four scalar lines attached to
it. For QED, the interaction
\[
    e\bar\psi\gamma^\mu\psi A_\mu
\]
has one photon line and two fermion lines attached to the vertex.

A Feynman diagram is therefore a picture of a term in the Wick expansion of the
Dyson series. It is not a literal photograph of particles moving in spacetime.

\section{External Lines, Internal Lines, and Vertices}

External lines represent field insertions associated with initial or final
states, or with operators whose correlation function is being computed. Internal
lines represent contractions between fields inside the perturbative expansion.
They carry propagators.

Vertices represent interaction terms. In \(\phi^4\) theory, the vertex factor is
related to \(-i\lambda\). In QED, the vertex factor is related to
\(-ie\gamma^\mu\), with conventions depending on the sign chosen for the
interaction Lagrangian.

In momentum space, each internal line carries a momentum variable. Momentum is
conserved at each vertex. Integrations over undetermined internal momenta become
loop integrals.

\section{Connected and Disconnected Diagrams}

A diagram is connected if every part of it is linked to every other part by
lines and vertices. A disconnected diagram factorizes into independent pieces.

In correlation functions, disconnected pieces often include vacuum bubbles,
which are diagrams with no external lines. The normalized generating formulas of
Chapter 7 divide by the vacuum persistence amplitude, causing pure vacuum
bubbles to cancel.

Connected diagrams are the ones that carry the essential correlation between
external insertions. Later, when we construct scattering amplitudes, connected
amputated diagrams will play the central role.

\section{Vacuum Bubbles}

Vacuum bubbles arise from contractions among fields belonging only to interaction
vertices, with no connection to external fields. For example, in \(\phi^4\)
theory, a first-order vacuum bubble comes from contracting the four fields at a
single interaction point among themselves.

Vacuum bubbles contribute to
\[
    \langle0|T\exp\left[-i\int\dd^4x\,\mathcal{H}_{\mathrm{int}}(x)\right]|0\rangle.
\]
They affect the vacuum persistence amplitude but cancel from normalized
correlation functions. This cancellation is why one often focuses on connected
diagrams attached to the external points.

\section{Symmetry Factors}

The same diagram can arise from many equivalent Wick contractions. The symmetry
factor accounts for this overcounting. In a perturbative expansion, the factors
\(1/n!\) from the Dyson series and factors such as \(1/4!\) in the interaction
Lagrangian combine with the number of contractions to produce the final
coefficient of a diagram.

A symmetry factor is the inverse of the number of graph automorphisms that leave
the diagram unchanged while preserving its external structure. In practice, one
often learns symmetry factors by examples.

For instance, in \(\phi^4\) theory, the factor \(1/4!\) is chosen so that the
basic four-point vertex has a simple Feynman rule. Without this factorial, every
vertex would carry extra combinatorial factors from permuting identical scalar
fields.

\section{Interpretation Limits of Feynman Diagrams}

Feynman diagrams are extraordinarily useful, but they should be interpreted
carefully. A diagram represents a term in a perturbative expansion, not a unique
classical history. Internal lines do not represent directly observed particles.
They represent propagators, or equivalently contractions of free fields.

This distinction matters especially for virtual particles. A virtual particle is
not a particle detected in an experiment. It is a feature of a diagrammatic
representation of a correlation function or amplitude. Different gauges or field
variables can reorganize the same physical amplitude into different sets of
intermediate diagrammatic terms.

The physical content lies in gauge-invariant amplitudes, cross sections, decay
rates, and correlation functions, not in any single diagram by itself.

\begin{summarybox}
Wick's theorem converts time-ordered products of free fields into sums of
normal-ordered products and contractions. Vacuum expectation values keep only
fully contracted terms, and each contraction is a propagator. Applying Wick's
theorem to the Dyson series produces Feynman diagrams: vertices come from
interactions, internal lines from propagators, external lines from field
insertions, and symmetry factors from combinatorics.
\end{summarybox}

\section{Chapter Checklist}

After reading this chapter, one should be able to:
\begin{itemize}
    \item explain why Wick's theorem is needed in perturbation theory;
    \item define normal ordering and contractions;
    \item state Wick's theorem for free scalar fields;
    \item evaluate a four-point free scalar vacuum expectation value;
    \item explain how contractions become propagators;
    \item state the fermionic version of Wick's theorem qualitatively;
    \item explain the origin of fermionic sign factors;
    \item translate Wick contractions into Feynman diagrams;
    \item distinguish external lines, internal lines, and vertices;
    \item explain connected diagrams, vacuum bubbles, and symmetry factors;
    \item state why Feynman diagrams should not be interpreted as literal
    particle trajectories.
\end{itemize}
